% TOP_PROBLEM: {"id": 30007355, "title": "Liouville Brownian motion as the walk limit on FK–Ising maps", "category_id": 19, "category": {"id": 19, "name": "probability", "display_name": "Probability", "description": "Problems involving probability theory, stochastic processes, and random structures.", "slug": "probability", "order_index": 19, "created_at": "2026-07-31T15:26:25.671Z"}, "status": "open"}
% FIRST_POSTED: 2026-10-06
% ENTRY_KIND: statement_only
% !TeX program = lualatex
\documentclass[11pt]{article}
\usepackage[margin=25mm]{geometry}
\usepackage{fontspec}
\setmainfont{Latin Modern Roman}
\usepackage{amsmath,amssymb,mathtools}
\usepackage{unicode-math}
\setmathfont{Latin Modern Math}
\usepackage{xurl,hyperref}
\hypersetup{colorlinks=true,urlcolor=blue}
\setcounter{secnumdepth}{0}
\setlength{\parindent}{0pt}
\setlength{\parskip}{5pt}
\setlength{\emergencystretch}{3em}
\newcommand{\R}{\mathbb R}
\newcommand{\N}{\mathbb N}
\newcommand{\Z}{\mathbb Z}
\newcommand{\Q}{\mathbb Q}
\newcommand{\C}{\mathbb C}
\newcommand{\D}{\mathbb D}
\newcommand{\F}{\mathbb F}
\newcommand{\sref}[2]{\hyperlink{source:#1:#2}{[S#2]}}
\newcommand{\cataloguescope}[1]{#1}
\begin{document}
\section*{Liouville Brownian motion as the walk limit on FK–Ising maps}\label{30007355}
\subsection{Definitions and mathematical statement}
\paragraph{Statement recorded in the supplied source.}
For the critical spherical FK–Ising n-edge map law, assume that a deterministic $a_n\to 0$ gives metric–measure convergence to the unit-area sqrt(3)-Liouville quantum sphere. Do any fixed finite number of conditionally independent simple random walks, each initially distributed according to the degree probability measure and accelerated by n, converge jointly with that space to independent stationary Liouville Brownian motions, in the curve-decorated Gromov–Hausdorff–Prokhorov topology?
\subsection{Short English statement}
For the critical spherical FK–Ising n-edge map law, assume that a deterministic $a_n\to 0$ gives metric–measure convergence to the unit-area sqrt(3)-Liouville quantum sphere. Do any fixed finite number of conditionally independent simple random walks, each initially distributed according to the degree probability measure and accelerated by n, converge jointly with that space to independent stationary Liouville Brownian motions, in the curve-decorated Gromov–Hausdorff–Prokhorov topology?
\subsection{Sources}
\begin{itemize}
\item[S1] Original question source or mathematical context.
\url{https://github.com/openai/math/blob/adc7f1241b42e322a6451854ab7e4b4c146bf78a/preprints/Random-Walks-on-Critical-FK-Ising-Maps-and-Liouville-Brownian-Motion-October-5-2026/build/sections/introduction.tex}.
\item[S2] OpenAI, manuscript citation page. Source of the formulation motivating this local problem.
\url{https://github.com/openai/math/blob/adc7f1241b42e322a6451854ab7e4b4c146bf78a/preprints/Random-Walks-on-Critical-FK-Ising-Maps-and-Liouville-Brownian-Motion-October-5-2026/README.md}.
\end{itemize}
\end{document}
